%% kapa: exam=1 round=27 year=2016 period=1 form=A subject=real_estate subject_ko=부동산학원론 seq=35 number=115
%% source: Q-Net 감정평가사 기출문제 (cst00309) artlSeq=5169402 file=extracted/1008/5169402/제27회 감정평가사 1차 문제/제27회 감정평가사 1교시문제(A형).pdf
%% conversion: kordoc 4.24.51 -> tools/gen1.py (auto); review=proofread-sonnet; flags=none
\begin{question}{115}
분양가상한제에 관한 설명으로 옳지 않은 것은?
\begin{choices}
  \item 주택구매 수요자들의 주택구입 부담을 덜어주기 위해 신규분양주택의 분양가격을 주택법령에 따라 정한 가격을 초과하여 받지 못하도록 규제하는 제도이다.
  \item 주택법령상 사업주체가 일반인에게 공급하는 공동주택 중 공공택지 외의 택지에서 주택가격 상승 우려가 있어 심의를 거쳐 지정하는 지역에서 공급하는 주택의 경우에는 기준에 따라 산정되는 분양가격 이하로 공급하여야 한다.
  \item 공급자의 이윤이 저하되어 주택의 공급이 감소하는 현상이 나타날 수 있다.
  \item 주택법령상 사업주체는 분양가상한제 적용주택으로서 공공택지에서 공급하는 주택에 대하여 입주자모집 승인을 받았을 때에는 입주자 모집공고에 택지비, 공사비, 간접비 등에 대하여 분양가격을 공시하여야 한다.
  \item 주택법령상 사업주체가 일반인에게 공급하는 공동주택 중 공공택지에서 공급하는 도시형 생활주택은 분양가상한제를 적용한다.
\end{choices}
\end{question}
\answer{5}
